\chapter{Code-Dokumentation}
\label{app:code}

\section{Überblick}

Dieser Anhang dokumentiert die wichtigsten Code-Komponenten der Transformer-Implementierung. Der vollständige Quellcode ist im GitHub-Repository verfügbar und umfasst etwa 2.500 Zeilen Python-Code mit ausführlicher Dokumentation.

\section{Projektstruktur}

\begin{lstlisting}[style=pythonstyle, caption=Projektstruktur]
src/
├── models/
│   ├── __init__.py
│   ├── transformer.py          # Vollständige Transformer-Implementierung
│   └── language_model.py       # GPT-style Language Model
├── utils/
│   ├── __init__.py
│   ├── tokenizer.py           # Einfacher Tokenizer
│   └── data_loader.py         # Datenverarbeitungsfunktionen
├── training/
│   ├── __init__.py
│   ├── trainer.py             # Training-Loop und Evaluierung
│   └── optimizer.py           # Optimierungsstrategien
└── experiments/
    ├── train_language_model.py # Beispiel-Training-Script
    ├── evaluate_model.py       # Evaluierungs-Script
    └── generate_text.py        # Textgenerierungs-Demo
\end{lstlisting}

\section{Kernkomponenten}

\subsection{Multi-Head Attention}

Die zentrale Attention-Implementierung:

\begin{lstlisting}[style=pythonstyle, caption=Multi-Head Attention Klasse]
class MultiHeadAttention(nn.Module):
    """
    Multi-Head Attention mechanism implementation.
    
    Args:
        d_model: Model dimension
        n_heads: Number of attention heads
        dropout: Dropout probability
    """
    
    def __init__(self, d_model: int, n_heads: int, dropout: float = 0.1):
        super().__init__()
        assert d_model % n_heads == 0
        
        self.d_model = d_model
        self.n_heads = n_heads
        self.d_k = d_model // n_heads
        
        # Linear projections for Q, K, V
        self.w_q = nn.Linear(d_model, d_model, bias=False)
        self.w_k = nn.Linear(d_model, d_model, bias=False)
        self.w_v = nn.Linear(d_model, d_model, bias=False)
        self.w_o = nn.Linear(d_model, d_model)
        
        self.dropout = nn.Dropout(dropout)
\end{lstlisting}

\subsection{Positional Encoding}

Sinusoidale Positional Encodings:

\begin{lstlisting}[style=pythonstyle, caption=Positional Encoding Implementation]
class PositionalEncoding(nn.Module):
    """
    Sinusoidal positional encoding implementation.
    """
    
    def __init__(self, d_model: int, max_len: int = 5000):
        super().__init__()
        
        pe = torch.zeros(max_len, d_model)
        position = torch.arange(0, max_len, dtype=torch.float).unsqueeze(1)
        
        div_term = torch.exp(
            torch.arange(0, d_model, 2).float() * 
            (-math.log(10000.0) / d_model)
        )
        
        pe[:, 0::2] = torch.sin(position * div_term)
        pe[:, 1::2] = torch.cos(position * div_term)
        pe = pe.unsqueeze(0).transpose(0, 1)
        
        self.register_buffer('pe', pe)
\end{lstlisting}

\subsection{Transformer Layer}

Encoder und Decoder Layer Implementierung:

\begin{lstlisting}[style=pythonstyle, caption=Encoder Layer]
class EncoderLayer(nn.Module):
    """
    Single Transformer Encoder Layer.
    """
    
    def __init__(self, d_model: int, n_heads: int, d_ff: int, dropout: float = 0.1):
        super().__init__()
        self.self_attention = MultiHeadAttention(d_model, n_heads, dropout)
        self.feed_forward = PositionwiseFeedForward(d_model, d_ff, dropout)
        self.norm1 = nn.LayerNorm(d_model)
        self.norm2 = nn.LayerNorm(d_model)
        self.dropout = nn.Dropout(dropout)
    
    def forward(self, x: torch.Tensor, mask: Optional[torch.Tensor] = None):
        # Self-attention with residual connection
        attention_output = self.self_attention(x, x, x, mask)
        x = self.norm1(x + self.dropout(attention_output))
        
        # Feed-forward with residual connection
        ff_output = self.feed_forward(x)
        x = self.norm2(x + self.dropout(ff_output))
        
        return x
\end{lstlisting}

\section{Training Infrastructure}

\subsection{Trainer Klasse}

Die Hauptklasse für Training und Evaluierung:

\begin{lstlisting}[style=pythonstyle, caption=Training Infrastructure]
class LanguageModelTrainer:
    """
    Comprehensive trainer for Transformer language models.
    
    Features:
    - Training with proper scheduling and optimization
    - Evaluation and monitoring
    - Checkpointing and model saving
    - Sample generation during training
    - Training curve visualization
    """
    
    def __init__(self, model, tokenizer, device, output_dir="./checkpoints"):
        self.model = model
        self.tokenizer = tokenizer
        self.device = device
        self.output_dir = output_dir
        
        # Training state
        self.global_step = 0
        self.epoch = 0
        self.best_loss = float('inf')
        
        # Logging
        self.train_losses = []
        self.eval_losses = []
        self.learning_rates = []
\end{lstlisting}

\subsection{Optimierung}

AdamW Optimizer mit Weight Decay:

\begin{lstlisting}[style=pythonstyle, caption=Optimizer Setup]
def create_optimizer(self, learning_rate=1e-4, weight_decay=0.01):
    """Create AdamW optimizer with proper weight decay."""
    
    # Separate parameters for weight decay
    decay_params = []
    no_decay_params = []
    
    for name, param in self.model.named_parameters():
        if param.requires_grad:
            if 'bias' in name or 'norm' in name:
                no_decay_params.append(param)
            else:
                decay_params.append(param)
    
    optimizer_grouped_parameters = [
        {'params': decay_params, 'weight_decay': weight_decay},
        {'params': no_decay_params, 'weight_decay': 0.0}
    ]
    
    return optim.AdamW(optimizer_grouped_parameters, lr=learning_rate)
\end{lstlisting}

\section{Tokenizer}

\subsection{Einfacher Tokenizer}

Word-level Tokenizer mit Vokabular-Management:

\begin{lstlisting}[style=pythonstyle, caption=Tokenizer Implementation]
class SimpleTokenizer:
    """
    Simple word-level tokenizer with vocabulary management.
    
    Features:
    - Text preprocessing and normalization
    - Vocabulary building from corpus
    - Encoding/decoding with special tokens
    - Batch processing with padding
    """
    
    def __init__(self, vocab_size: int = 10000, min_freq: int = 2):
        self.vocab_size = vocab_size
        self.min_freq = min_freq
        
        # Special tokens
        self.special_tokens = {
            '<pad>': 0, '<unk>': 1, '<bos>': 2, '<eos>': 3
        }
        
        self.token_to_id = {}
        self.id_to_token = {}
        self.token_freq = {}
    
    def build_vocabulary(self, texts: List[str]):
        """Build vocabulary from training texts."""
        token_counter = Counter()
        
        for text in texts:
            tokens = self.tokenize(text)
            token_counter.update(tokens)
        
        # Filter by frequency and build vocab
        # ... implementation details
\end{lstlisting}

\section{Modell-Konfiguration}

\subsection{Factory Functions}

Einfache Erstellung verschiedener Modell-Konfigurationen:

\begin{lstlisting}[style=pythonstyle, caption=Model Factory Functions]
def create_small_language_model(vocab_size: int, **kwargs) -> GPTLanguageModel:
    """Create a small language model for experimentation."""
    return GPTLanguageModel(
        vocab_size=vocab_size,
        d_model=kwargs.get('d_model', 256),
        n_heads=kwargs.get('n_heads', 4),
        n_layers=kwargs.get('n_layers', 4),
        d_ff=kwargs.get('d_ff', 1024),
        max_len=kwargs.get('max_len', 512),
        dropout=kwargs.get('dropout', 0.1)
    )

def create_medium_language_model(vocab_size: int, **kwargs) -> GPTLanguageModel:
    """Create a medium-sized language model."""
    return GPTLanguageModel(
        vocab_size=vocab_size,
        d_model=kwargs.get('d_model', 384),
        n_heads=kwargs.get('n_heads', 6),
        n_layers=kwargs.get('n_layers', 6),
        d_ff=kwargs.get('d_ff', 1536),
        max_len=kwargs.get('max_len', 512),
        dropout=kwargs.get('dropout', 0.1)
    )
\end{lstlisting}

\section{Beispiel-Scripts}

\subsection{Training Script}

Vollständiges Trainings-Beispiel:

\begin{lstlisting}[style=pythonstyle, caption=Training Example]
def main():
    # Setup
    device = torch.device("cuda" if torch.cuda.is_available() else "cpu")
    
    # Load data and create tokenizer
    texts = load_sample_data()
    tokenizer = SimpleTokenizer(vocab_size=5000)
    tokenizer.build_vocabulary(texts)
    
    # Create model
    model = create_small_language_model(
        vocab_size=tokenizer.get_vocab_size()
    ).to(device)
    
    # Prepare data
    dataset = create_simple_dataset(texts, tokenizer, max_length=128)
    train_dataloader, eval_dataloader = create_data_loaders(
        dataset, batch_size=32
    )
    
    # Train
    trainer = LanguageModelTrainer(model, tokenizer, device)
    trainer.train(
        train_dataloader=train_dataloader,
        eval_dataloader=eval_dataloader,
        num_epochs=10,
        learning_rate=1e-4
    )
\end{lstlisting}

\subsection{Textgenerierung}

Interaktive Textgenerierung:

\begin{lstlisting}[style=pythonstyle, caption=Text Generation Script]
@torch.no_grad()
def generate_text_interactive(model, tokenizer, device):
    """Interactive text generation loop."""
    model.eval()
    
    print("Interactive Text Generation (type 'quit' to exit)")
    
    while True:
        prompt = input("\nEnter prompt: ")
        if prompt.lower() == 'quit':
            break
        
        # Encode prompt
        input_ids = torch.tensor(
            tokenizer.encode(prompt), dtype=torch.long
        ).unsqueeze(0).to(device)
        
        # Generate
        generated = model.generate(
            input_ids=input_ids,
            max_new_tokens=50,
            temperature=0.8,
            top_k=50,
            do_sample=True
        )
        
        # Decode and display
        generated_text = tokenizer.decode(generated[0].cpu().tolist())
        print(f"Generated: {generated_text}")
\end{lstlisting}

\section{Testing und Validierung}

\subsection{Unit Tests}

Beispiel für Komponenten-Tests:

\begin{lstlisting}[style=pythonstyle, caption=Unit Test Example]
import pytest
import torch
from src.models.transformer import MultiHeadAttention

def test_multihead_attention_output_shape():
    """Test that MultiHeadAttention produces correct output shape."""
    batch_size, seq_len, d_model = 2, 10, 512
    n_heads = 8
    
    attention = MultiHeadAttention(d_model, n_heads)
    
    # Create test input
    x = torch.randn(batch_size, seq_len, d_model)
    
    # Forward pass
    output = attention(x, x, x)
    
    # Check output shape
    assert output.shape == (batch_size, seq_len, d_model)

def test_attention_with_mask():
    """Test attention mechanism with masking."""
    # ... test implementation
\end{lstlisting}

\section{Performance Monitoring}

\subsection{Metriken und Logging}

Umfassendes Monitoring während Training:

\begin{lstlisting}[style=pythonstyle, caption=Monitoring Implementation]
def log_training_metrics(self, loss, learning_rate, gradient_norm):
    """Log comprehensive training metrics."""
    
    metrics = {
        'epoch': self.epoch,
        'global_step': self.global_step,
        'train_loss': loss,
        'learning_rate': learning_rate,
        'gradient_norm': gradient_norm,
        'perplexity': math.exp(loss),
        'tokens_per_second': self.calculate_throughput(),
        'gpu_memory_used': torch.cuda.memory_allocated() / 1e9
    }
    
    # Log to various backends
    self.log_to_tensorboard(metrics)
    self.log_to_wandb(metrics)
    self.save_metrics_json(metrics)

def calculate_throughput(self):
    """Calculate tokens processed per second."""
    if hasattr(self, '_last_time') and hasattr(self, '_last_tokens'):
        current_time = time.time()
        current_tokens = self.global_step * self.batch_size * self.seq_len
        
        time_diff = current_time - self._last_time
        token_diff = current_tokens - self._last_tokens
        
        return token_diff / time_diff if time_diff > 0 else 0
    
    self._last_time = time.time()
    self._last_tokens = 0
    return 0
\end{lstlisting}

\section{Konfiguration und Reproduzierbarkeit}

\subsection{Hyperparameter Management}

Strukturierte Konfigurationsverwaltung:

\begin{lstlisting}[style=pythonstyle, caption=Configuration Management]
from dataclasses import dataclass
from typing import Optional

@dataclass
class ModelConfig:
    """Model architecture configuration."""
    vocab_size: int
    d_model: int = 512
    n_heads: int = 8
    n_layers: int = 6
    d_ff: int = 2048
    max_len: int = 512
    dropout: float = 0.1

@dataclass
class TrainingConfig:
    """Training configuration."""
    batch_size: int = 32
    learning_rate: float = 1e-4
    weight_decay: float = 0.01
    warmup_steps: int = 1000
    max_epochs: int = 10
    gradient_clip_norm: float = 1.0
    save_steps: int = 1000
    eval_steps: int = 500

def set_random_seeds(seed: int = 42):
    """Set all random seeds for reproducibility."""
    import random
    import numpy as np
    
    random.seed(seed)
    np.random.seed(seed)
    torch.manual_seed(seed)
    torch.cuda.manual_seed(seed)
    torch.cuda.manual_seed_all(seed)
    
    # For deterministic behavior (may impact performance)
    torch.backends.cudnn.deterministic = True
    torch.backends.cudnn.benchmark = False
\end{lstlisting}

\section{Zusammenfassung}

Die Code-Implementierung umfasst:

\begin{itemize}
    \item \textbf{Vollständige Transformer-Architektur} mit allen wichtigen Komponenten
    \item \textbf{Modulare Struktur} für einfache Erweiterung und Wartung
    \item \textbf{Umfassende Training-Infrastructure} mit Monitoring und Checkpointing
    \item \textbf{Flexible Konfiguration} für verschiedene Experimente
    \item \textbf{Ausführliche Dokumentation} und Beispiele
    \item \textbf{Testing-Framework} für Validierung der Komponenten
    \item \textbf{Reproduzierbarkeit} durch Seed-Kontrolle und Konfiguration
\end{itemize}

Der vollständige Code ist im GitHub-Repository verfügbar und kann als Ausgangspunkt für weitere Experimente und Entwicklungen verwendet werden.